% Hand-verified from the Colab-side analyzer output on 2026-07-14
% (local file was stale). Colab-side numbers reflect 225 real HEM rows
% (H1 + M3 at N=25 × 3 paths, all strict-AND-valid closures).
%
% Regenerate by syncing results/results_holdout_humaneval_mutated.tsv
% from Drive and re-running analyze/holdout_sensitivity.py.
\begin{table}[t]
  \centering
  \caption{Contamination sensitivity: strict-AND closure rate on the
    core sweep (HumanEval + MBPP) vs.\ the HumanEval-Mutated held-out
    subset (HEM-25: function-rename + docstring-paraphrase over the
    same underlying algorithms; $n = 25$ per cell per path), for the
    hetero champion H1 and the strongest mono baseline M3. Large
    negative $\Delta$ on HEM-25 would signal surface-form
    memorisation; the observed pattern (positive $\Delta$ on the
    behavioural paths, negative $\Delta$ on the surface-form BERTScore
    path) instead indicates that the core closure rates are not driven
    by contamination.
    \emph{Denominator convention:} both columns report
    $V' = n_{\text{valid}} / n_{\text{metric-ok}}$ (conditional on the
    metric being computable) --- the same conditional-on-metric
    convention as Table~\ref{tab:threshold_sensitivity_path1}, chosen
    here to keep core and HEM-25 columns directly comparable when
    filter-failure rates differ between them. Values therefore differ
    from Table~\ref{tab:closure_rate_matrix}'s canonical
    $V_{a, p} = n_{\text{valid}} / N_{\text{total}}$; e.g.\ H1 Path~1
    is $64/81 = 79.0\%$ here vs.\ $64/89 = 71.9\%$ under canonical
    $V$. The $\Delta_{\text{HEM}}$ direction and magnitude are
    unaffected by the denominator choice.
    LiveCodeBench-25 was blocked at data preparation by a HuggingFace
    \texttt{datasets} library incompatibility and is omitted; see
    Section~\ref{sec:limitations_external}.}
  \label{tab:holdout_sensitivity}
  \begin{tabular}{llrrr}
    \toprule
    Cell & Path & Core & HEM-25 & $\Delta_{\textsf{HEM}}$ \\
    \midrule
    H1 & 1 (mutation kill rate)    & $79.0\%$          & $100.0\%$         & $+21.0$~pp \\
    H1 & 2 (reference-test pass)   & $64.0\%$          & $\phantom{0}79.6\%$ & $+15.5$~pp \\
    H1 & 3 (BERTScore F1)          & $63.6\%$          & $\phantom{0}32.7\%$ & $\mathbf{-31.0}$~\textbf{pp} \\
    \midrule
    M3 & 1 (mutation kill rate)    & $85.5\%$          & $100.0\%$         & $+14.5$~pp \\
    M3 & 2 (reference-test pass)   & $61.2\%$          & $\phantom{0}87.0\%$ & $+25.7$~pp \\
    M3 & 3 (BERTScore F1)          & $70.0\%$          & $\phantom{0}45.5\%$ & $\mathbf{-24.5}$~\textbf{pp} \\
    \bottomrule
  \end{tabular}

  \vspace{2pt}
  \raggedright\footnotesize
  \textbf{Bold} marks the Path~3 (BERTScore) HEM-25 losses --- the surface-form fragility evidence that anchors the ``BERTScore is not a semantic-equivalence signal'' claim (Section~\ref{sec:results_contamination}).
\end{table}
